\documentclass{article}

\usepackage{geometry}
\geometry{letterpaper}

% style options: numeric, authoryear, alphabetic,
%                ieee, apa, chicago-authordate, ...
\usepackage[backend=biber, style=apa]{biblatex}
\addbibresource{logic1.bib}

\usepackage{amssymb}
\usepackage{outlines}

\usepackage{enumitem}
\setenumerate[1]{label=\Roman*.}
\setenumerate[2]{label=\Alph*.}
\setenumerate[3]{label=\arabic*.}

\title{Formal Logic 1}
\author{Christopher R. Walker}
\date{\today}

\begin{document}

\maketitle

\begin{outline}[enumerate]

\1 \textbf{Week 1: What Is Logic?}
\\
\\
\textbf{Reading:}
\cite[][Chapters 1-3 (pp. 2-25)]{Mag25}
\\
\\
\textbf{Focus:}
validity, soundness, entailment, satisfiability;
distinguishing deductive from inductive arguments

\2 Arguments
\2 The scope of logic
\2 Other logical notions

\1 \textbf{Week 2: Symbolizing English}
\\
\\
\textbf{Reading:}
\cite[][Chapters 4-5 (pp. 27-48)]{Mag25}
\\
\\
\textbf{Focus:}
translating natural language into TFL;
understanding connective meanings

\2 First steps to symbolization
\2 Connectives

\1 \textbf{Week 3: TFL Syntax}
\\
\\
\textbf{Reading:}
\cite[][Chapters 6-8 (pp. 49-67)]{Mag25}
\\
\\
\textbf{Focus:}
well-formed formulas, parse trees, scope, avoiding ambiguity

\2 Sentences of TFL
\2 Ambiguity
\2 Use and mention

\1 \textbf{Week 4: TFL Semantics}
\\
\\
\textbf{Reading:}
\cite[][Chapters 9-11 (pp. 69-84)]{Mag25}
\\
\\
\textbf{Focus:}
understanding truth-functional connectives;
constructing truth tables for complex sentences;

\2 Characteristic truth tables
\2 Truth-functional connectives
\2 Complete truth tables

\1 \textbf{Week 5: Semantic Concepts}
\\
\\
\textbf{Reading:}
\cite[][Chapters 12-14 (pp. 85-101)]{Mag25}
\\
\\
\textbf{Focus:}
using truth tables to test logical properties;
efficiency strategies

\2 Semantic concepts
\2 Limitations of TFL
\2 Truth table shortcuts

\1 \textbf{Week 6: Partial Truth Tables}
\\
\\
\textbf{Reading:}
\cite[][Chapter 15 (pp. 102-8)]{Mag25}
\\
\\
\textbf{Focus:}
using partial tables for validation;
catch-up and review

\2 Partial Truth Tables

\1 \textbf{Week 7: Introduction to Natural Deduction}
\\
\\
\textbf{Focus:}
the Fitch-style proof system;
introduction and elimination rules
\2 The very idea of natural deduction \3 \cite{Mag25},  Chapter 16 (pp 110-2)
\2 Basic rules for TFL \3 \cite{Mag25},  Chapter 17 (pp 113-41)

\1 \textbf{Week 8: Constructing Proofs}
\\
\\
\textbf{Focus:}
proof strategies;
working backward from the conclusion;
subproofs
\2 Constructing proofs \3 \cite{Mag25},  Chapter 18 (pp 142-161)

\1 \textbf{Week 9: Additional Rules \& Proof-Theoretic Concepts}
\\
\\
\textbf{Focus:}
extending the basic system;
connecting proof theory to semantics
\2 Additional rules for TFL \3 \cite{Mag25},  Chapter 19 (pp 162-9)
\2 Proof-theoretic concepts \3 \cite{Mag25},  Chapter 20 (pp 170-3)
\2 Derived rules \3 \cite{Mag25},  Chapter 21 (pp 174-81)

\1 \textbf{Week 10: Soundness, Completeness}
\\
\\
\textbf{Focus:}
relating semantic and proof-theoretic consequence;
reviewing TFL natural deduction
\2 Soundness and completeness \3 \cite{Mag25},  Chapter 22 (pp 182-8)

\1 \textbf{Week 11: Building Blocks of FOL}
\\
\\
\textbf{Focus:}
translating simple universal and existential claims;
predicate notation
\2 Building Blocks of FOL \3 \cite{Mag25},  Chapter 23 (pp 190-8)
\2 Sentences with One Quantifier \3 \cite{Mag25},  Chapter 24 (pp 199-212)

\1 \textbf{Week 12: Multiple Generality \& Identity}
\\
\\
\textbf{Focus:}
nested quantifiers;
the identity predicate;
definite descriptions
\2 Multiple generality \3 \cite{Mag25},  Chapter 25 (pp 213-27)
\2 Identity \3 \cite{Mag25},  Chapter 26 (pp 228-34)
\2 Sentences of FOL \3 \cite{Mag25},  Chapter 27 (pp235-42)

\1 \textbf{Week 13: Definite Descriptions \& FOL Ambiguity}
\\
\\
\textbf{Focus:}
Russellian analysis of "the";
scope ambiguities in FOL
\2 Definite descriptions \3 \cite{Mag25},  Chapter 28 (pp 243-52)
\2 Ambiguity \3 \cite{Mag25},  Chapter 29 (pp 253-6)

\1 \textbf{Week 14: FOL Semantics}
\\
\\
\textbf{Focus:}
defining truth relative to an interpretation;
testing validity via counter-models
\2 Extensionality \3 \cite{Mag25},  Chapter 30 (pp 258-66)
\2 Truth in FOL \3 \cite{Mag25},  Chapter 31 (pp 267-76)
\2 Semantic concepts \3 \cite{Mag25},  Chapter 32 (pp 277-80)
\2 Using interpretations \3 \cite{Mag25},  Chapter 33 (pp 281-8)
\2 Reasoning about interpretations \3 \cite{Mag25},  Chapter 34 (pp 289-92)
\2 ** Properties of relations \3 \cite{Mag25},  Chapter 35 (pp 293-5)

\1 \textbf{Week 15: FOL Proofs}
\\
\\
\textbf{Focus:}
universal and existential introduction/elimination;
identity rules;
connecting FOL proof theory to semantics
\2 Basic rules for FOL \3 \cite{Mag25},  Chapter 36 (pp 297-315)
\2 Proofs with quantifiers \3 \cite{Mag25},  Chapter 37 (pp 316-22)
\2 Conversion of quantifiers \3 \cite{Mag25},  Chapter 38 (pp 323-5)
\2 Rules for identity \3 \cite{Mag25},  Chapter 39 (pp 326-9)
\2 ** Derived rules \3 \cite{Mag25},  Chapter 40 (pp 330-1)
\2 ** Proofs and semantics \3 \cite{Mag25},  Chapter 41 (pp 332-5)

\end{outline}

\printbibliography

\end{document}  
